%%%%%%%%%%%%%%%%%%%%%%%%%%%%%%%%%%%%%%%%%%%%%%%%%%%%%%%%%%%%%%%%%%%%%%%%%%%%%%%%
% ACF-DOC-039.02
% Severe Weather Prediction System
%%%%%%%%%%%%%%%%%%%%%%%%%%%%%%%%%%%%%%%%%%%%%%%%%%%%%%%%%%%%%%%%%%%%%%%%%%%%%%%%

\section{Severe Weather Prediction System}

\subsection{Executive Summary}

The Severe Weather Prediction System is a core component of the
Atmospheric Complexity Framework (ACF) Multi-Hazard Decision Support
platform.

Its objective is to detect, predict and monitor high-impact atmospheric
phenomena using a combination of:

\begin{itemize}

\item Numerical Weather Prediction models

\item Radar observations

\item Satellite observations

\item Surface observations

\item Atmospheric soundings

\item Data assimilation

\item Artificial intelligence

\item Earth System Digital Twin

\end{itemize}

The system focuses on hazardous convective phenomena including:

\begin{itemize}

\item Severe thunderstorms

\item Supercells

\item Mesoscale Convective Systems

\item Tornadoes

\item Hailstorms

\item Lightning events

\item Downbursts

\item Extreme convective precipitation

\end{itemize}

The final objective is to provide probabilistic forecasts and
operational warnings through the Atmospheric Weather Center Interface
(AWCI).

%%%%%%%%%%%%%%%%%%%%%%%%%%%%%%%%%%%%%%%%%%%%%%%%%%%%%%%%%%%%%%%%%%%%%%%%%%%%%%%%

\subsection{Introduction}

Severe weather events are among the most complex atmospheric phenomena
because they involve interactions between thermodynamics, dynamics,
microphysics and surface processes.

Unlike large-scale weather systems, severe convection develops on
spatial and temporal scales that require high-resolution observations
and numerical simulations.

Typical severe storms involve:

\begin{itemize}

\item Strong atmospheric instability

\item High moisture availability

\item Dynamic lifting mechanisms

\item Significant vertical wind shear

\item Organized cloud structures

\end{itemize}

The ACF system represents severe weather as a coupled atmospheric
process inside the Earth System Digital Twin.

%%%%%%%%%%%%%%%%%%%%%%%%%%%%%%%%%%%%%%%%%%%%%%%%%%%%%%%%%%%%%%%%%%%%%%%%%%%%%%%%

\subsection{Atmospheric Convection Fundamentals}

Convection occurs when an air parcel becomes warmer and less dense than
its surrounding environment.

The vertical acceleration of an air parcel can be expressed as:

\[
\frac{dw}{dt}=B-g
\]

where:

\begin{itemize}

\item $w$ is vertical velocity

\item $B$ is buoyancy acceleration

\item $g$ is gravitational acceleration

\end{itemize}


The buoyancy term is approximated by:

\[
B=g
\frac{T_v'-T_v}{T_v}
\]

where:

\begin{itemize}

\item $T_v'$ is parcel virtual temperature

\item $T_v$ is environmental virtual temperature

\end{itemize}

Positive buoyancy promotes upward motion and cloud development.

%%%%%%%%%%%%%%%%%%%%%%%%%%%%%%%%%%%%%%%%%%%%%%%%%%%%%%%%%%%%%%%%%%%%%%%%%%%%%%%%

\subsection{Convective Storm Ingredients}

The development of severe convection requires four main ingredients.

\subsubsection{Instability}

Instability represents the potential energy available for convection.

The most common parameter is Convective Available Potential Energy:

\[
CAPE=
\int_{LFC}^{EL}
g
\frac{T_v'-T_v}{T_v}
dz
\]

where:

\begin{itemize}

\item LFC = Level of Free Convection

\item EL = Equilibrium Level

\end{itemize}


Large CAPE values indicate stronger potential updrafts.

%%%%%%%%%%%%%%%%%%%%%%%%%%%%%%%%%%%%%%%%%%%%%%%%%%%%%%%%%%%%%%%%%%%%%%%%%%%%%%%%

\subsubsection{Moisture}

Atmospheric moisture controls cloud formation and precipitation.

Important ACF parameters include:

\begin{itemize}

\item Specific humidity

\item Mixing ratio

\item Relative humidity

\item Dew point temperature

\item Precipitable water

\end{itemize}

%%%%%%%%%%%%%%%%%%%%%%%%%%%%%%%%%%%%%%%%%%%%%%%%%%%%%%%%%%%%%%%%%%%%%%%%%%%%%%%%

\subsubsection{Lift}

A lifting mechanism initiates convection.

ACF analyses:

\begin{itemize}

\item Cold fronts

\item Warm fronts

\item Dry lines

\item Sea breeze boundaries

\item Orographic lifting

\item Outflow boundaries

\end{itemize}

%%%%%%%%%%%%%%%%%%%%%%%%%%%%%%%%%%%%%%%%%%%%%%%%%%%%%%%%%%%%%%%%%%%%%%%%%%%%%%%%

\subsubsection{Wind Shear}

Vertical wind shear controls storm organization.

The bulk shear vector is:

\[
\vec{S}
=
\vec{V}_{top}
-
\vec{V}_{bottom}
\]

Strong shear favors organized severe storms such as supercells.

%%%%%%%%%%%%%%%%%%%%%%%%%%%%%%%%%%%%%%%%%%%%%%%%%%%%%%%%%%%%%%%%%%%%%%%%%%%%%%%%
%%%%%%%%%%%%%%%%%%%%%%%%%%%%%%%%%%%%%%%%%%%%%%%%%%%%%%%%%%%%%%%%%%%%%%%%%%%%%%%%
\subsection{Severe Thunderstorm Classification}

The ACF Severe Weather Engine classifies convective systems according to
their structure, dynamics and associated hazards.

The main categories are:

\begin{itemize}

\item Ordinary thunderstorms

\item Multicell thunderstorms

\item Supercell thunderstorms

\item Mesoscale Convective Systems (MCS)

\item Squall lines

\item Derecho events

\end{itemize}

The classification combines:

\begin{itemize}

\item Radar reflectivity

\item Doppler velocity

\item Satellite cloud properties

\item Numerical model output

\item Artificial intelligence classification

\end{itemize}

%%%%%%%%%%%%%%%%%%%%%%%%%%%%%%%%%%%%%%%%%%%%%%%%%%%%%%%%%%%%%%%%%%%%%%%%%%%%%%%%

\subsection{Ordinary and Multicell Thunderstorms}

Ordinary thunderstorms are short-lived convective cells generally
controlled by local instability.

Their life cycle contains:

\begin{enumerate}

\item Developing stage

\item Mature stage

\item Dissipating stage

\end{enumerate}


Multicell storms consist of several convective cells at different stages
of development.

They can produce:

\begin{itemize}

\item Heavy precipitation

\item Lightning

\item Localized flooding

\item Strong wind gusts

\end{itemize}

%%%%%%%%%%%%%%%%%%%%%%%%%%%%%%%%%%%%%%%%%%%%%%%%%%%%%%%%%%%%%%%%%%%%%%%%%%%%%%%%

\subsection{Supercell Dynamics}

Supercells represent one of the most dangerous forms of convection.

They are characterized by:

\begin{itemize}

\item Persistent rotating updraft

\item Mesocyclone formation

\item Strong vertical wind shear

\item Long lifetime

\item Severe weather production

\end{itemize}


The rotating updraft is produced by the interaction between horizontal
vorticity and vertical motion.

The vertical vorticity equation can be approximated as:

\[
\frac{\partial \zeta}{\partial t}
=
-\vec{V}\cdot\nabla\zeta
+
Stretching
+
Tilting
+
Friction
\]

where:

\begin{itemize}

\item $\zeta$ represents vertical vorticity

\item Stretching increases rotation intensity

\item Tilting converts horizontal rotation into vertical rotation

\end{itemize}

%%%%%%%%%%%%%%%%%%%%%%%%%%%%%%%%%%%%%%%%%%%%%%%%%%%%%%%%%%%%%%%%%%%%%%%%%%%%%%%%

\subsection{Mesocyclone Detection}

The ACF radar intelligence module detects mesocyclones using:

\begin{itemize}

\item Doppler velocity couplets

\item Rotation signatures

\item Vertical continuity

\item Reflectivity structures

\end{itemize}


The rotation parameter is calculated as:

\[
Rotation =
\frac{\Delta V}{\Delta x}
\]

where:

\begin{itemize}

\item $\Delta V$ is radial velocity difference

\item $\Delta x$ is horizontal distance

\end{itemize}

%%%%%%%%%%%%%%%%%%%%%%%%%%%%%%%%%%%%%%%%%%%%%%%%%%%%%%%%%%%%%%%%%%%%%%%%%%%%%%%%

\subsection{Mesoscale Convective Systems}

Mesoscale Convective Systems (MCS) are organized groups of thunderstorms
covering large areas.

Typical characteristics:

\begin{itemize}

\item Horizontal scale greater than 100 km

\item Long duration

\item Organized precipitation structures

\item Strong wind hazards

\end{itemize}


MCS categories include:

\begin{itemize}

\item Mesoscale Convective Complexes

\item Squall Lines

\item Bow Echoes

\item Derechos

\end{itemize}

%%%%%%%%%%%%%%%%%%%%%%%%%%%%%%%%%%%%%%%%%%%%%%%%%%%%%%%%%%%%%%%%%%%%%%%%%%%%%%%%

\subsection{Squall Lines}

Squall lines are linear convective systems associated with strong
downdrafts and damaging winds.

Their structure contains:

\begin{itemize}

\item Leading convective line

\item Cold pool

\item Rear inflow jet

\item Stratiform precipitation region

\end{itemize}


The ACF system monitors:

\begin{itemize}

\item Line orientation

\item Propagation speed

\item Wind intensity

\item Lightning density

\end{itemize}

%%%%%%%%%%%%%%%%%%%%%%%%%%%%%%%%%%%%%%%%%%%%%%%%%%%%%%%%%%%%%%%%%%%%%%%%%%%%%%%%

\subsection{Derecho Events}

Derechos are widespread severe wind events generated by organized
convective systems.

Operational indicators include:

\begin{itemize}

\item Bow echo structures

\item Strong rear inflow jet

\item Long-lived convection

\item Extensive damage path

\end{itemize}


ACF integrates:

\begin{itemize}

\item Radar analysis

\item Wind field prediction

\item AI pattern recognition

\item Impact assessment

\end{itemize}

%%%%%%%%%%%%%%%%%%%%%%%%%%%%%%%%%%%%%%%%%%%%%%%%%%%%%%%%%%%%%%%%%%%%%%%%%%%%%%%%

\subsection{Tornado Environment Prediction}

Tornado prediction requires analysis of:

\begin{itemize}

\item Instability

\item Low-level moisture

\item Vertical wind shear

\item Storm-relative helicity

\item Mesocyclone potential

\end{itemize}


The Storm Relative Helicity (SRH) is:

\[
SRH =
-\int
(\vec{V}-\vec{C})
\cdot
\nabla\times\vec{V}
dz
\]

where:

\begin{itemize}

\item $\vec{V}$ is wind vector

\item $\vec{C}$ is storm motion vector

\end{itemize}

%%%%%%%%%%%%%%%%%%%%%%%%%%%%%%%%%%%%%%%%%%%%%%%%%%%%%%%%%%%%%%%%%%%%%%%%%%%%%%%%
%%%%%%%%%%%%%%%%%%%%%%%%%%%%%%%%%%%%%%%%%%%%%%%%%%%%%%%%%%%%%%%%%%%%%%%%%%%%%%%%
\subsection{Advanced Severe Weather Parameters}

The ACF Severe Weather Prediction Engine uses a large set of
thermodynamic and dynamic parameters to diagnose the potential for
hazardous convection.

These parameters are calculated from atmospheric soundings, numerical
weather prediction models and assimilated observations.

%%%%%%%%%%%%%%%%%%%%%%%%%%%%%%%%%%%%%%%%%%%%%%%%%%%%%%%%%%%%%%%%%%%%%%%%%%%%%%%%

\subsection{Convective Inhibition (CIN)}

Convective Inhibition represents the amount of energy required to lift
an air parcel until it reaches free convection.

It is defined as:

\[
CIN =
-\int_{Surface}^{LFC}
g
\frac{T_v'-T_v}{T_v}
dz
\]

Large CIN values indicate a stable atmosphere and suppressed convection.

Small CIN values combined with strong instability indicate an environment
favorable for explosive storm development.

%%%%%%%%%%%%%%%%%%%%%%%%%%%%%%%%%%%%%%%%%%%%%%%%%%%%%%%%%%%%%%%%%%%%%%%%%%%%%%%%

\subsection{Lifted Index (LI)}

The Lifted Index measures atmospheric stability.

It is defined as:

\[
LI=T_{500}-T_{p500}
\]

where:

\begin{itemize}

\item $T_{500}$ is environmental temperature at 500 hPa

\item $T_{p500}$ is parcel temperature at 500 hPa

\end{itemize}


Operational interpretation:

\begin{center}

\begin{tabular}{lll}

\toprule

LI Value & Stability & Risk \\

\midrule

Positive & Stable & Low convection \\

0 to -3 & Weak instability & Possible storms \\

-3 to -6 & Moderate instability & Severe storms possible \\

Below -6 & Strong instability & Extreme convection potential \\

\bottomrule

\end{tabular}

\end{center}

%%%%%%%%%%%%%%%%%%%%%%%%%%%%%%%%%%%%%%%%%%%%%%%%%%%%%%%%%%%%%%%%%%%%%%%%%%%%%%%%

\subsection{Total Totals Index}

The Total Totals index combines thermal instability and low-level
moisture.

\[
TT =
(T_{850}-T_{500})
+
(Td_{850}-T_{500})
\]

Higher values indicate increasing thunderstorm potential.

%%%%%%%%%%%%%%%%%%%%%%%%%%%%%%%%%%%%%%%%%%%%%%%%%%%%%%%%%%%%%%%%%%%%%%%%%%%%%%%%

\subsection{K Index}

The K Index estimates thunderstorm potential using temperature and
humidity profiles.

\[
K =
(T_{850}-T_{500})
+
Td_{850}
-
(T_{700}-Td_{700})
\]

It evaluates:

\begin{itemize}

\item Low-level moisture

\item Mid-level lapse rates

\item Atmospheric dryness

\end{itemize}

%%%%%%%%%%%%%%%%%%%%%%%%%%%%%%%%%%%%%%%%%%%%%%%%%%%%%%%%%%%%%%%%%%%%%%%%%%%%%%%%

\subsection{SWEAT Index}

The Severe Weather Threat Index combines several ingredients:

\[
SWEAT =
12Td_{850}
+
20(TT-49)
+
2V_{850}
+
V_{500}
+
125(Shear)
\]

It includes:

\begin{itemize}

\item Moisture

\item Instability

\item Wind speed

\item Wind shear

\end{itemize}

%%%%%%%%%%%%%%%%%%%%%%%%%%%%%%%%%%%%%%%%%%%%%%%%%%%%%%%%%%%%%%%%%%%%%%%%%%%%%%%%

\subsection{Bulk Richardson Number (BRN)}

The Bulk Richardson Number compares buoyancy and wind shear.

\[
BRN=
\frac{CAPE}
{0.5(V_{shear})^2}
\]

Interpretation:

\begin{itemize}

\item High BRN: pulse convection

\item Moderate BRN: organized storms

\item Low BRN: supercell environment

\end{itemize}

%%%%%%%%%%%%%%%%%%%%%%%%%%%%%%%%%%%%%%%%%%%%%%%%%%%%%%%%%%%%%%%%%%%%%%%%%%%%%%%%

\subsection{Storm Relative Helicity (SRH)}

Storm Relative Helicity measures the potential for rotating updrafts.

\[
SRH=
-\int_{0}^{z}
(\vec{V}-\vec{C})
\cdot
\nabla\times\vec{V}
dz
\]

Important layers:

\begin{itemize}

\item 0--1 km SRH

\item 0--3 km SRH

\item Effective SRH

\end{itemize}

High SRH values increase tornado potential when combined with
instability.

%%%%%%%%%%%%%%%%%%%%%%%%%%%%%%%%%%%%%%%%%%%%%%%%%%%%%%%%%%%%%%%%%%%%%%%%%%%%%%%%

\subsection{Significant Tornado Parameter (STP)}

The Significant Tornado Parameter combines multiple tornado ingredients.

A simplified formulation:

\[
STP=
\frac{CAPE}{1500}
\times
\frac{SRH}{150}
\times
\frac{Shear}{20}
\times
\frac{LCL}{1000}
\]

The parameter combines:

\begin{itemize}

\item Instability

\item Low-level rotation

\item Deep-layer shear

\item Cloud base height

\end{itemize}

%%%%%%%%%%%%%%%%%%%%%%%%%%%%%%%%%%%%%%%%%%%%%%%%%%%%%%%%%%%%%%%%%%%%%%%%%%%%%%%%

\subsection{Supercell Composite Parameter (SCP)}

The Supercell Composite Parameter estimates supercell potential.

\[
SCP=
f(CAPE)
\times
f(SRH)
\times
f(Shear)
\]

High values indicate environments favorable for rotating storms.

%%%%%%%%%%%%%%%%%%%%%%%%%%%%%%%%%%%%%%%%%%%%%%%%%%%%%%%%%%%%%%%%%%%%%%%%%%%%%%%%

\subsection{LCL, LFC and Equilibrium Level}

The ACF system calculates important vertical atmospheric levels.

\subsubsection{Lifted Condensation Level}

LCL represents cloud base formation.

\[
LCL=f(T,Td,P)
\]


\subsubsection{Level of Free Convection}

LFC represents the altitude where a parcel becomes positively buoyant.

\[
LFC=z(B>0)
\]


\subsubsection{Equilibrium Level}

EL represents the top of the convective cloud.

\[
EL=z(B=0)
\]

%%%%%%%%%%%%%%%%%%%%%%%%%%%%%%%%%%%%%%%%%%%%%%%%%%%%%%%%%%%%%%%%%%%%%%%%%%%%%%%%

\subsection{Hodograph Analysis}

The hodograph represents wind evolution with altitude.

It provides information about:

\begin{itemize}

\item Wind shear

\item Storm motion

\item Rotation potential

\item Supercell characteristics

\end{itemize}


The ACF visualization engine displays:

\begin{itemize}

\item Surface to 6 km wind profile

\item SRH regions

\item Storm-relative flow

\item Convective environment

\end{itemize}

%%%%%%%%%%%%%%%%%%%%%%%%%%%%%%%%%%%%%%%%%%%%%%%%%%%%%%%%%%%%%%%%%%%%%%%%%%%%%%%%
%%%%%%%%%%%%%%%%%%%%%%%%%%%%%%%%%%%%%%%%%%%%%%%%%%%%%%%%%%%%%%%%%%%%%%%%%%%%%%%%
\subsection{Radar Observation Integration}

Weather radar represents one of the most important observation systems
for severe weather monitoring.

The ACF Severe Weather Engine integrates multi-source radar information:

\begin{itemize}

\item Reflectivity

\item Doppler radial velocity

\item Dual polarization products

\item Echo tops

\item Vertically Integrated Liquid (VIL)

\item Hydrometeor classification

\end{itemize}


%%%%%%%%%%%%%%%%%%%%%%%%%%%%%%%%%%%%%%%%%%%%%%%%%%%%%%%%%%%%%%%%%%%%%%%%%%%%%%%%

\subsubsection{Radar Reflectivity}

Radar reflectivity provides information about precipitation intensity.

The reflectivity factor is:

\[
Z=
\sum D^6
\]

where:

\begin{itemize}

\item $Z$ represents radar reflectivity

\item $D$ represents hydrometeor diameter

\end{itemize}


The ACF system analyzes:

\begin{itemize}

\item High reflectivity cores

\item Hail signatures

\item Hook echoes

\item Bow echoes

\item Weak echo regions

\end{itemize}


%%%%%%%%%%%%%%%%%%%%%%%%%%%%%%%%%%%%%%%%%%%%%%%%%%%%%%%%%%%%%%%%%%%%%%%%%%%%%%%%

\subsubsection{Doppler Velocity Analysis}

Doppler radar measures radial wind velocity.

The measured velocity is:

\[
V_r =
V
\cos(\theta)
\]

where:

\begin{itemize}

\item $V_r$ is radial velocity

\item $V$ is wind velocity

\item $\theta$ is observation angle

\end{itemize}


The ACF engine detects:

\begin{itemize}

\item Rotation signatures

\item Mesocyclones

\item Tornadic vortex signatures

\item Strong wind zones

\end{itemize}


%%%%%%%%%%%%%%%%%%%%%%%%%%%%%%%%%%%%%%%%%%%%%%%%%%%%%%%%%%%%%%%%%%%%%%%%%%%%%%%%

\subsection{Satellite Observation Integration}

Satellite observations provide continuous global monitoring.

The ACF platform integrates:

\begin{itemize}

\item Geostationary satellites

\item Polar orbiting satellites

\item Infrared imagery

\item Visible imagery

\item Water vapor channels

\item Lightning mapping observations

\end{itemize}


Important satellite indicators:

\begin{itemize}

\item Overshooting tops

\item Rapid cloud growth

\item Cold cloud tops

\item Enhanced water vapor divergence

\item Convective initiation zones

\end{itemize}


%%%%%%%%%%%%%%%%%%%%%%%%%%%%%%%%%%%%%%%%%%%%%%%%%%%%%%%%%%%%%%%%%%%%%%%%%%%%%%%%

\subsection{Artificial Intelligence Severe Weather Engine}

The ACF AI Severe Weather Engine combines physical prediction and
machine learning techniques.

The hybrid concept is:

\[
Prediction =
NWP
+
Observation
+
AI
\]


The AI system performs:

\begin{itemize}

\item Storm classification

\item Object detection

\item Tracking

\item Intensity estimation

\item Probability forecasting

\end{itemize}


%%%%%%%%%%%%%%%%%%%%%%%%%%%%%%%%%%%%%%%%%%%%%%%%%%%%%%%%%%%%%%%%%%%%%%%%%%%%%%%%

\subsection{Deep Learning Architecture}

The AI architecture includes:

\begin{itemize}

\item Convolutional Neural Networks for radar images

\item Transformers for temporal evolution

\item Graph Neural Networks for atmospheric interactions

\item Physics-Informed Neural Networks

\end{itemize}


The input data include:

\begin{itemize}

\item Radar sequences

\item Satellite images

\item Numerical model fields

\item Soundings

\item Surface observations

\end{itemize}


The output provides:

\begin{itemize}

\item Storm probability

\item Hazard category

\item Intensity prediction

\item Forecast confidence

\end{itemize}


%%%%%%%%%%%%%%%%%%%%%%%%%%%%%%%%%%%%%%%%%%%%%%%%%%%%%%%%%%%%%%%%%%%%%%%%%%%%%%%%

\subsection{Storm Tracking System}

The ACF tracking engine follows severe weather objects in time.

Tracked parameters:

\begin{itemize}

\item Position

\item Movement vector

\item Propagation speed

\item Intensity evolution

\item Hazard probability

\end{itemize}


The storm motion vector is:

\[
\vec{C}
=
(C_x,C_y)
\]


Future position estimation:

\[
X(t+\Delta t)
=
X(t)+C\Delta t
\]


%%%%%%%%%%%%%%%%%%%%%%%%%%%%%%%%%%%%%%%%%%%%%%%%%%%%%%%%%%%%%%%%%%%%%%%%%%%%%%%%

\subsection{Integration with Earth System Digital Twin}

The Severe Weather Engine continuously exchanges information with the
ACF Earth System Digital Twin.

The Digital Twin provides:

\begin{itemize}

\item Atmospheric state

\item Moisture distribution

\item Wind fields

\item Temperature structure

\item Surface conditions

\item Ocean influence

\end{itemize}


The Digital Twin enables:

\begin{itemize}

\item Real-time storm evolution

\item Scenario simulation

\item Probabilistic forecasting

\item Impact assessment

\end{itemize}


%%%%%%%%%%%%%%%%%%%%%%%%%%%%%%%%%%%%%%%%%%%%%%%%%%%%%%%%%%%%%%%%%%%%%%%%%%%%%%%%

\subsection{AWCI Operational Integration}

The Atmospheric Weather Center Interface receives severe weather
information from the prediction engine.

The AWCI dashboard displays:

\begin{itemize}

\item Severe weather maps

\item Radar overlays

\item Satellite layers

\item Storm tracks

\item Hazard probabilities

\item Warning levels

\end{itemize}


Operational workflow:

\[
Detection
\rightarrow
Analysis
\rightarrow
Forecast
\rightarrow
Warning
\rightarrow
Decision
\]


%%%%%%%%%%%%%%%%%%%%%%%%%%%%%%%%%%%%%%%%%%%%%%%%%%%%%%%%%%%%%%%%%%%%%%%%%%%%%%%%
%%%%%%%%%%%%%%%%%%%%%%%%%%%%%%%%%%%%%%%%%%%%%%%%%%%%%%%%%%%%%%%%%%%%%%%%%%%%%%%%
\subsection{Operational Severe Weather Workflow}

The ACF Severe Weather Prediction System operates continuously inside
the global operational environment.

The complete workflow is:

\[
Observation
\rightarrow
Quality Control
\rightarrow
Assimilation
\rightarrow
Prediction
\rightarrow
Detection
\rightarrow
Risk Assessment
\rightarrow
Warning
\]

The operational cycle includes:

\begin{enumerate}

\item Acquisition of global observations

\item Real-time atmospheric analysis

\item High-resolution numerical prediction

\item AI-based severe weather detection

\item Expert analysis support

\item Warning generation

\item Event monitoring

\end{enumerate}


%%%%%%%%%%%%%%%%%%%%%%%%%%%%%%%%%%%%%%%%%%%%%%%%%%%%%%%%%%%%%%%%%%%%%%%%%%%%%%%%

\subsection{High Performance Computing Requirements}

Severe weather prediction requires large computational resources because
of the high spatial and temporal resolution required.

The ACF HPC architecture provides:

\begin{itemize}

\item Parallel numerical simulation

\item GPU acceleration

\item Distributed data processing

\item AI model inference

\item Real-time visualization

\end{itemize}


The computational workflow is:

\[
Data
\rightarrow
HPC
\rightarrow
Model
\rightarrow
AI
\rightarrow
Visualization
\]


%%%%%%%%%%%%%%%%%%%%%%%%%%%%%%%%%%%%%%%%%%%%%%%%%%%%%%%%%%%%%%%%%%%%%%%%%%%%%%%%

\subsection{Forecast Verification}

The ACF system evaluates prediction quality using multiple verification
methods.

Verification parameters include:

\begin{itemize}

\item Probability of Detection (POD)

\item False Alarm Ratio (FAR)

\item Critical Success Index (CSI)

\item Bias Score

\item Forecast uncertainty

\end{itemize}


The Critical Success Index is:

\[
CSI=
\frac{Hits}
{Hits+Misses+False\ Alarms}
\]


%%%%%%%%%%%%%%%%%%%%%%%%%%%%%%%%%%%%%%%%%%%%%%%%%%%%%%%%%%%%%%%%%%%%%%%%%%%%%%%%

\subsection{Uncertainty Estimation}

Severe weather prediction contains significant uncertainty.

ACF uses:

\begin{itemize}

\item Ensemble Prediction Systems

\item Probabilistic forecasting

\item AI uncertainty estimation

\item Physics-based sensitivity analysis

\end{itemize}


The final prediction includes:

\[
Forecast =
Probability
+
Confidence
+
Uncertainty
\]


%%%%%%%%%%%%%%%%%%%%%%%%%%%%%%%%%%%%%%%%%%%%%%%%%%%%%%%%%%%%%%%%%%%%%%%%%%%%%%%%

\subsection{Human Expert Interaction}

Although ACF provides advanced automation, human expertise remains
essential.

The system supports forecasters through:

\begin{itemize}

\item Interactive analysis tools

\item AI explanations

\item Situation summaries

\item Risk prioritization

\item Decision assistance

\end{itemize}


%%%%%%%%%%%%%%%%%%%%%%%%%%%%%%%%%%%%%%%%%%%%%%%%%%%%%%%%%%%%%%%%%%%%%%%%%%%%%%%%

\subsection{Integration with Emergency Management}

Severe weather information is transmitted to the ACF emergency
management system.

The system provides:

\begin{itemize}

\item Geographic warning zones

\item Population exposure estimation

\item Infrastructure risk assessment

\item Response recommendations

\end{itemize}


%%%%%%%%%%%%%%%%%%%%%%%%%%%%%%%%%%%%%%%%%%%%%%%%%%%%%%%%%%%%%%%%%%%%%%%%%%%%%%%%

\subsection{Future Development}

Future versions of the ACF Severe Weather Engine will introduce:

\begin{itemize}

\item Kilometer-scale global convection prediction

\item Autonomous AI forecasting agents

\item Real-time urban hazard simulation

\item Advanced radar artificial intelligence

\item Fully coupled atmosphere-ocean-land prediction

\item Global severe weather digital twin

\end{itemize}


%%%%%%%%%%%%%%%%%%%%%%%%%%%%%%%%%%%%%%%%%%%%%%%%%%%%%%%%%%%%%%%%%%%%%%%%%%%%%%%%

\subsection{References}

Scientific foundations are based on:

\begin{itemize}

\item Atmospheric convection theory

\item Numerical Weather Prediction

\item Radar meteorology

\item Satellite meteorology

\item Mesoscale meteorology

\item Machine learning for weather prediction

\end{itemize}


%%%%%%%%%%%%%%%%%%%%%%%%%%%%%%%%%%%%%%%%%%%%%%%%%%%%%%%%%%%%%%%%%%%%%%%%%%%%%%%%

